
intéressant que modulo torsion, i.e. tensorisé par $\Q$. On voit en même temps, à l'aide de cet exemple, que ce n'est que tensorisé par $\Q$ que $\operatorname{Gr}^\bullet(X)$ à un caractère covariant par rapport aux morphismes propres, même en se bornant aux immersions fermées de schémas quasi-projectifs et lisses. En effet, si $X$ est tel que $j$ ne soit pas un isomorphisme, il existe dans $X$ un sous-schéma fermé intègre $Y$ tel que $\mathrm{c}\ell^\bullet(Y)\in K^\bullet(X)$ ne soit pas dans $\operatorname{Filt}^i(X)$, où $i=\operatorname{codim}(Y,X)$. Nous pouvons prendre $Y$ de dimension minimale, de sorte que $\operatorname{Filt}^{i+1}(X)=\operatorname{Filt}_{\mathrm{top}}^{i+1}(X)$. Si $Z$ est l'ensemble singulier de $Y$, et $X'=X-Z$, $Y'=Y-Z$, on en conclut que $\mathrm{c}\ell^\bullet(Y')\in K^\bullet(X')$ n'est pas dans $\operatorname{Filt}^i(X')$ (compte tenu que $K^\bullet(X)\to K^\bullet(X')$ est surjectif et compatible avec les augmentations, et que son noyau est $\operatorname{Im}\bigl(K(Z)\longrightarrow K(X)\bigr)\subset\operatorname{Filt}_{\mathrm{top}}^{i+1}(X)=\operatorname{Filt}_{\mathrm{top}}^{i+1}(X)$). Considérons alors l'inclusion
\[
 f:Y'\longrightarrow X',
\]
qui est une immersion fermée de codimension $i$ de schémas quasi-projectifs et lisses sur $k$, \uline{telle que}
\[
 f_*:K^\bullet(Y')\longrightarrow K^\bullet(X')
\]
\uline{ne soit pas compatible avec les} $\lambda$-\uline{filtrations} (\uline{modulo décalage par} $i$), puisque $f_*(1)$ n'est pas dans $\operatorname{Filt}^i(X')$.

\textbf{4.7.} Gardant les notations de 4.5, notons maintenant qu'il résulte de (Exp. IX) que le théorème de HODGKIN cité équivaut au fait que l'on a un isomorphisme
\[
 K^\bullet(G)\simeq\Z,\qquad\text{donc}\qquad \Gtop^\bullet(G)\simeq\Z,
\]
l'homomorphisme $K^\bullet(G/B)\to K^\bullet(G)$ étant à priori surjectif et ayant comme noyau l'idéal engendré par les $\mathrm{c}\ell^\bullet(\underline{L})-1$, où $\underline{L}$ parcourt les Modules inversibles sur $X=G/B$ associés aux représentations $B\to\underline{G}_m$ (qui donnent en fait tous les modules inversibles sur $X$ [5, Exp. 15]). On voit d'autre part par un raisonnement analogue que
\[
 A(G/B)\longrightarrow A(G)
\]
est également surjectif, et que son noyau est l'idéal engendré par les éléments de $A^1(G/B)$ définis par les représentations $B\to\underline{G}_m$, donc en fait par $A^1(G/B)$ tout entier.\SGAAuthorFootnote{*}{Cf. IX 3.4, 2.2.} Cela montre que $A(G)$ est réduit à $A^0(X)=\Z$ si et
